% Generated draft: requires mathematical and source review.
\documentclass[11pt]{article}
\usepackage[margin=25mm]{geometry}
\usepackage{fontspec}
\setmainfont{Latin Modern Roman}
\usepackage{amsmath,amssymb,mathtools,unicode-math}
\setmathfont{Latin Modern Math}
\usepackage{enumitem,needspace,xurl,hyperref}
\setlength{\parindent}{0pt}
\setlength{\parskip}{5pt}
\setlength{\emergencystretch}{3em}
\newcommand{\N}{\mathbb{N}}
\newcommand{\Z}{\mathbb{Z}}
\newcommand{\Q}{\mathbb{Q}}
\newcommand{\R}{\mathbb{R}}
\newcommand{\C}{\mathbb{C}}
\newcommand{\F}{\mathbb{F}}
\newcommand{\A}{\mathbb{A}}
\newcommand{\PP}{\mathbb{P}}
\newcommand{\E}{\mathbb{E}}
\newcommand{\eps}{\varepsilon}
\newcommand{\dd}{\,\mathrm{d}}
\begin{document}
\section*{UnsolvedMath ID: 701317}
\textbf{Source problem code:} catalogue entry
\subsection{Definitions and mathematical statement}
Let \(R_n=\Z[x_1,\ldots,x_n]\) be the polynomial ring in the finitely many indeterminates appearing in the input specification. Under the interpretation that \(\operatorname{Aut}(R_n)\) denotes the group of unital ring automorphisms, these automorphisms fix \(\Z\) pointwise. For \(f\in R_n\), define its stabilizer by
\[
G_f=\{\sigma\in\operatorname{Aut}(R_n):\sigma(f)=f\}.
\]
This is the unique largest subgroup all of whose elements fix \(f\). Here the source's phrase ``maximal subgroup'' is interpreted in this sense; it does not impose the additional condition that \(G_f\) be a maximal proper subgroup of \(\operatorname{Aut}(R_n)\).

\textbf{Requested algorithm.} Construct an algorithm which, given \(f\in\Z[x_1,\ldots,x_n]\), terminates and explicitly describes generators \(S_f\subseteq\operatorname{Aut}(R_n)\) such that
\[
\langle S_f\rangle=G_f.
\]
An explicit individual generator can be specified by its polynomial images of \(x_1,\ldots,x_n\). The source does not specify whether the required output is a finite list of generators or whether an effectively described infinite generating family is permitted. Consequently, the required representation of \(S_f\), and hence the precise algorithmic output requirement, remains unresolved.

The source relates this question to describing the integer solution set
\[
V_{\Z}(f)=\{(a_1,\ldots,a_n)\in\Z^n:f(a_1,\ldots,a_n)=0\}.
\]
It does not state a separate requirement that the algorithm describe this solution set.
\subsection{Short English statement}
Given an integer polynomial, find explicit generators for the group of polynomial ring automorphisms that fix it. The source leaves unclear whether finitely many generators must be returned or an effective description of an infinite family is allowed. The supplied registry labels the problem open; that label has not been independently verified.
\subsection{Sources}
\begin{itemize}
\item Supplied local UnsolvedMath registry record: permanent ID \(701317\), exact source code \texttt{catalogue entry}, title \emph{Kourovka Notebook Problem 17.15}. The record attributes the problem to V. G. Bardakov and gives the proposal year as 2010.
\item \emph{The Kourovka Notebook (Unsolved Problems in Group Theory)}, 17th issue (2010), Problem 17.15; also identified in the supplied background with the 21st issue edition (2026), \texttt{arXiv:1401.0300v48}. The supplied collection description names E. I. Khukhro and V. D. Mazurov as editors. These bibliographic details are supplied attribution and have not been independently checked.
\item The registry records status \texttt{open} and an update timestamp of \texttt{2026-10-16T14:00:00Z}. This is supplied registry metadata, not an independently verified literature assessment or a stated literature-verification date.
\end{itemize}
\end{document}
